\documentclass[11pt]{article}
\usepackage{amsmath, amssymb}
\usepackage{geometry}
\geometry{margin=1in}
\title{\textbf{APM1220: Applied Multivariate Data Analysis}\\ Formative Assessment 1}
\author{Aldrin A. Tupas}
\date{August 27, 2026}

\begin{document}

\maketitle
% ---------------------------------------------------------
% PROBLEM 2.2
% ---------------------------------------------------------
\section*{Problem 2.2}
Given the matrices
\[
A = \begin{bmatrix} -1 & 3 \\ 4 & 2 \end{bmatrix}, \quad
B = \begin{bmatrix} 4 & -3 \\ 1 & -2 \\ -2 & 0 \end{bmatrix}, \quad \text{and} \quad
C = \begin{bmatrix} 5 \\ -4 \\ 2 \end{bmatrix}
\]

\subsection*{(a) $5A$}
\[
5A = 5 \begin{bmatrix} -1 & 3 \\ 4 & 2 \end{bmatrix} 
= \begin{bmatrix} 5(-1) & 5(3) \\ 5(4) & 5(2) \end{bmatrix} 
= \mathbf{\begin{bmatrix} -5 & 15 \\ 20 & 10 \end{bmatrix}}
\]

\subsection*{(b) $BA$}
Matrix $B$ has dimension $3 \times 2$ and $A$ has dimension $2 \times 2$. Since the inner dimensions match, $BA$ is defined ($3 \times 2$):
\[
BA = \begin{bmatrix} 4 & -3 \\ 1 & -2 \\ -2 & 0 \end{bmatrix} \begin{bmatrix} -1 & 3 \\ 4 & 2 \end{bmatrix}
= \begin{bmatrix} 4(-1) + (-3)(4) & 4(3) + (-3)(2) \\ 1(-1) + (-2)(4) & 1(3) + (-2)(2) \\ -2(-1) + 0(4) & -2(3) + 0(2) \end{bmatrix}
= \mathbf{\begin{bmatrix} -16 & 6 \\ -9 & -1 \\ 2 & -6 \end{bmatrix}}
\]

\subsection*{(c) $A'B'$}
By transpose properties, $A'B' = (BA)'$. Taking the transpose of $BA$ from part (b):
\[
A'B' = \begin{bmatrix} -16 & 6 \\ -9 & -1 \\ 2 & -6 \end{bmatrix}' 
= \mathbf{\begin{bmatrix} -16 & -9 & 2 \\ 6 & -1 & -6 \end{bmatrix}}
\]

\subsection*{(d) $C'B$}
$C'$ has dimension $1 \times 3$ and $B$ has dimension $3 \times 2$. The multiplication yields a $1 \times 2$ row vector:
\[
C' = \begin{bmatrix} 5 & -4 & 2 \end{bmatrix}
\]
\[
C'B = \begin{bmatrix} 5 & -4 & 2 \end{bmatrix} \begin{bmatrix} 4 & -3 \\ 1 & -2 \\ -2 & 0 \end{bmatrix}
= \begin{bmatrix} 5(4) + (-4)(1) + 2(-2) & 5(-3) + (-4)(-2) + 2(0) \end{bmatrix}
\]
\[
= \begin{bmatrix} 20 - 4 - 4 & -15 + 8 + 0 \end{bmatrix} 
= \mathbf{\begin{bmatrix} 12 & -7 \end{bmatrix}}
\]

\subsection*{(e) Is $AB$ defined?}
$A$ is of dimension $2 \times 2$ and $B$ is of dimension $3 \times 2$. 
Since the number of columns of $A$ ($2$) does not equal the number of rows of $B$ ($3$), \textbf{$AB$ is not defined}.

% ---------------------------------------------------------
% PROBLEM 2.4
% ---------------------------------------------------------
\section*{Problem 2.4}

\subsection*{(a) Prove $(A')^{-1} = (A^{-1})'$}
\textbf{Proof:}
By definition of the inverse matrix, we know that:
\[
A A^{-1} = I
\]
Taking the transpose of both sides:
\[
(A A^{-1})' = I'
\]
Using the property of transpose for products $(XY)' = Y'X'$, and noting that $I' = I$:
\[
(A^{-1})' A' = I
\]
Similarly, starting from $A^{-1}A = I$:
\[
(A^{-1}A)' = I' \implies A' (A^{-1})' = I
\]
By definition, if $X Y = Y X = I$, then $Y = X^{-1}$. Letting $X = A'$ and $Y = (A^{-1})'$, we conclude:
\[
\mathbf{(A')^{-1} = (A^{-1})'}
\]
\hfill $\blacksquare$

\subsection*{(b) Prove $(AB)^{-1} = B^{-1}A^{-1}$}
\textbf{Proof:}
By definition, a matrix $M$ is the inverse of $(AB)$ if $(AB)M = I$ and $M(AB) = I$. 

Let us test $M = B^{-1}A^{-1}$:
1. Right multiplication:
\[
(AB)(B^{-1}A^{-1}) = A(B B^{-1})A^{-1} = A I A^{-1} = A A^{-1} = I
\]
2. Left multiplication:
\[
(B^{-1}A^{-1})(AB) = B^{-1}(A^{-1}A)B = B^{-1} I B = B^{-1}B = I
\]
Since both products yield the identity matrix $I$, $B^{-1}A^{-1}$ is the unique inverse of $AB$:
\[
\mathbf{(AB)^{-1} = B^{-1}A^{-1}}
\]
\hfill $\blacksquare$

% ---------------------------------------------------------
% PROBLEM 2.8
% ---------------------------------------------------------
\section*{Problem 2.8}
Given the matrix
\[
A = \begin{bmatrix} 1 & 2 \\ 2 & -2 \end{bmatrix}
\]

\subsection*{1. Eigenvalues $\lambda_1, \lambda_2$}
Solve the characteristic equation $\det(A - \lambda I) = 0$:
\[
\det \begin{bmatrix} 1-\lambda & 2 \\ 2 & -2-\lambda \end{bmatrix} = (1-\lambda)(-2-\lambda) - (2)(2) = 0
\]
\[
\lambda^2 + \lambda - 2 - 4 = 0 \implies \lambda^2 + \lambda - 6 = 0
\]
\[
(\lambda + 3)(\lambda - 2) = 0
\]
Thus, the eigenvalues are \textbf{$\lambda_1 = 2$} and \textbf{$\lambda_2 = -3$}.

\subsection*{2. Normalized Eigenvectors $e_1, e_2$}

\textbf{For $\lambda_1 = 2$:}
Solve $(A - 2I)x = 0$:
\[
\begin{bmatrix} -1 & 2 \\ 2 & -4 \end{bmatrix} \begin{bmatrix} x_1 \\ x_2 \end{bmatrix} = \begin{bmatrix} 0 \\ 0 \end{bmatrix} \implies -x_1 + 2x_2 = 0 \implies x_1 = 2x_2
\]
An eigenvector is $x^{(1)} = \begin{bmatrix} 2 \\ 1 \end{bmatrix}$. Normalizing gives:
\[
\|x^{(1)}\| = \sqrt{2^2 + 1^2} = \sqrt{5} \implies \mathbf{e_1 = \begin{bmatrix} \frac{2}{\sqrt{5}} \\ \frac{1}{\sqrt{5}} \end{bmatrix}}
\]

\textbf{For $\lambda_2 = -3$:}
Solve $(A - (-3)I)x = (A + 3I)x = 0$:
\[
\begin{bmatrix} 4 & 2 \\ 2 & 1 \end{bmatrix} \begin{bmatrix} x_1 \\ x_2 \end{bmatrix} = \begin{bmatrix} 0 \\ 0 \end{bmatrix} \implies 2x_1 + x_2 = 0 \implies x_2 = -2x_1
\]
An eigenvector is $x^{(2)} = \begin{bmatrix} 1 \\ -2 \end{bmatrix}$. Normalizing gives:
\[
\|x^{(2)}\| = \sqrt{1^2 + (-2)^2} = \sqrt{5} \implies \mathbf{e_2 = \begin{bmatrix} \frac{1}{\sqrt{5}} \\ -\frac{2}{\sqrt{5}} \end{bmatrix}}
\]

\subsection*{3. Spectral Decomposition}
The spectral decomposition of $A$ is $A = \lambda_1 e_1 e_1' + \lambda_2 e_2 e_2'$:
\[
e_1 e_1' = \begin{bmatrix} \frac{2}{\sqrt{5}} \\ \frac{1}{\sqrt{5}} \end{bmatrix} \begin{bmatrix} \frac{2}{\sqrt{5}} & \frac{1}{\sqrt{5}} \end{bmatrix} = \begin{bmatrix} \frac{4}{5} & \frac{2}{5} \\[4pt] \frac{2}{5} & \frac{1}{5} \end{bmatrix}
\]
\[
e_2 e_2' = \begin{bmatrix} \frac{1}{\sqrt{5}} \\ -\frac{2}{\sqrt{5}} \end{bmatrix} \begin{bmatrix} \frac{1}{\sqrt{5}} & -\frac{2}{\sqrt{5}} \end{bmatrix} = \begin{bmatrix} \frac{1}{5} & -\frac{2}{5} \\[4pt] -\frac{2}{5} & \frac{4}{5} \end{bmatrix}
\]
Thus, the spectral decomposition is:
\[
\mathbf{A = 2 \begin{bmatrix} \frac{4}{5} & \frac{2}{5} \\[4pt] \frac{2}{5} & \frac{1}{5} \end{bmatrix} - 3 \begin{bmatrix} \frac{1}{5} & -\frac{2}{5} \\[4pt] -\frac{2}{5} & \frac{4}{5} \end{bmatrix}}
\]

% ---------------------------------------------------------
% PROBLEM 2.11
% ---------------------------------------------------------
\section*{Problem 2.11}
\textbf{Proof by Induction:}
Let $A$ be a $p \times p$ diagonal matrix:
\[
A = \begin{bmatrix} 
a_{11} & 0 & 0 & \cdots & 0 \\
0 & a_{22} & 0 & \cdots & 0 \\
0 & 0 & a_{33} & \cdots & 0 \\
\vdots & \vdots & \vdots & \ddots & \vdots \\
0 & 0 & 0 & \cdots & a_{pp}
\end{bmatrix}
\]
Expanding the determinant along the first row:
\[
|A| = \sum_{j=1}^p a_{1j} A_{1j}
\]
where $A_{1j}$ is the cofactor. Since $a_{1j} = 0$ for all $j \neq 1$, this simplifies to:
\[
|A| = a_{11} A_{11} + 0 + \dots + 0 = a_{11} |A^{(1)}|
\]
where $A^{(1)}$ is the $(p-1) \times (p-1)$ diagonal submatrix obtained by deleting the first row and first column of $A$:
\[
A^{(1)} = \begin{bmatrix} 
a_{22} & 0 & \cdots & 0 \\
0 & a_{33} & \cdots & 0 \\
\vdots & \vdots & \ddots & \vdots \\
0 & 0 & \cdots & a_{pp}
\end{bmatrix}
\]
Applying cofactor expansion along the first row of $A^{(1)}$:
\[
|A^{(1)}| = a_{22} |A^{(2)}|
\]
Repeating this process iteratively for $k = 1, 2, \dots, p-1$ yields:
\[
|A| = a_{11} \cdot a_{22} \cdot |A^{(2)}| = a_{11} \cdot a_{22} \cdots a_{pp}
\]
Thus, the determinant of a diagonal matrix is the product of its diagonal elements:
\[
\mathbf{|A| = \prod_{i=1}^p a_{ii} = a_{11}a_{22}\cdots a_{pp}}
\]
\hfill $\blacksquare$

% ---------------------------------------------------------
% PROBLEM 2.24
% ---------------------------------------------------------
\section*{Problem 2.24}
Given $\Sigma = \begin{bmatrix} 4 & 0 & 0 \\ 0 & 9 & 0 \\ 0 & 0 & 1 \end{bmatrix}$.

\subsection*{(a) Find $\Sigma^{-1}$}
For a diagonal matrix, the inverse is obtained by taking the reciprocals of the diagonal elements:
\[
\mathbf{\Sigma^{-1} = \begin{bmatrix} \frac{1}{4} & 0 & 0 \\[4pt] 0 & \frac{1}{9} & 0 \\[4pt] 0 & 0 & 1 \end{bmatrix}}
\]

\subsection*{(b) Eigenvalues and Eigenvectors of $\Sigma$}
Since $\Sigma$ is diagonal, its eigenvalues are the diagonal entries:
\[
\mathbf{\lambda_1 = 9, \quad \lambda_2 = 4, \quad \lambda_3 = 1}
\]
The corresponding normalized eigenvectors are the standard basis vectors:
\[
\mathbf{e_1 = \begin{bmatrix} 0 \\ 1 \\ 0 \end{bmatrix}, \quad e_2 = \begin{bmatrix} 1 \\ 0 \\ 0 \end{bmatrix}, \quad e_3 = \begin{bmatrix} 0 \\ 0 \\ 1 \end{bmatrix}}
\]

\subsection*{(c) Eigenvalues and Eigenvectors of $\Sigma^{-1}$}
The eigenvectors of $\Sigma^{-1}$ are identical to those of $\Sigma$, and its eigenvalues are the reciprocals of the eigenvalues of $\Sigma$:
\[
\mathbf{\lambda_1^* = \frac{1}{9}, \quad \lambda_2^* = \frac{1}{4}, \quad \lambda_3^* = 1}
\]
The corresponding normalized eigenvectors are:
\[
\mathbf{e_1^* = \begin{bmatrix} 0 \\ 1 \\ 0 \end{bmatrix}, \quad e_2^* = \begin{bmatrix} 1 \\ 0 \\ 0 \end{bmatrix}, \quad e_3^* = \begin{bmatrix} 0 \\ 0 \\ 1 \end{bmatrix}}
\]

% ---------------------------------------------------------
% PROBLEM 2.30
% ---------------------------------------------------------
\section*{Problem 2.30}
Given $\mu_X = \begin{bmatrix} 4 \\ 3 \\ 2 \\ 1 \end{bmatrix}$, $\Sigma_X = \begin{bmatrix} 3 & 0 & 2 & 2 \\ 0 & 1 & 1 & 0 \\ 2 & 1 & 9 & -2 \\ 2 & 0 & -2 & 4 \end{bmatrix}$, partitioned as:
\[
X^{(1)} = \begin{bmatrix} X_1 \\ X_2 \end{bmatrix}, \quad X^{(2)} = \begin{bmatrix} X_3 \\ X_4 \end{bmatrix}
\]
\[
\mu^{(1)} = \begin{bmatrix} 4 \\ 3 \end{bmatrix}, \quad \mu^{(2)} = \begin{bmatrix} 2 \\ 1 \end{bmatrix}
\]
\[
\Sigma_{11} = \begin{bmatrix} 3 & 0 \\ 0 & 1 \end{bmatrix}, \quad \Sigma_{22} = \begin{bmatrix} 9 & -2 \\ -2 & 4 \end{bmatrix}, \quad \Sigma_{12} = \begin{bmatrix} 2 & 2 \\ 1 & 0 \end{bmatrix}, \quad \Sigma_{21} = \begin{bmatrix} 2 & 1 \\ 2 & 0 \end{bmatrix}
\]
with $A = \begin{bmatrix} 1 & 2 \end{bmatrix}$ and $B = \begin{bmatrix} 1 & -2 \\ 2 & -1 \end{bmatrix}$.

\subsection*{(a) $E(X^{(1)})$}
\[
\mathbf{E(X^{(1)}) = \mu^{(1)} = \begin{bmatrix} 4 \\ 3 \end{bmatrix}}
\]

\subsection*{(b) $E(AX^{(1)})$}
\[
E(AX^{(1)}) = A E(X^{(1)}) = \begin{bmatrix} 1 & 2 \end{bmatrix} \begin{bmatrix} 4 \\ 3 \end{bmatrix} = 1(4) + 2(3) = \mathbf{10}
\]

\subsection*{(c) $Cov(X^{(1)})$}
\[
\mathbf{Cov(X^{(1)}) = \Sigma_{11} = \begin{bmatrix} 3 & 0 \\ 0 & 1 \end{bmatrix}}
\]

\subsection*{(d) $Cov(AX^{(1)})$}
\[
Cov(AX^{(1)}) = A \Sigma_{11} A' = \begin{bmatrix} 1 & 2 \end{bmatrix} \begin{bmatrix} 3 & 0 \\ 0 & 1 \end{bmatrix} \begin{bmatrix} 1 \\ 2 \end{bmatrix}
= \begin{bmatrix} 3 & 2 \end{bmatrix} \begin{bmatrix} 1 \\ 2 \end{bmatrix} = 3(1) + 2(2) = \mathbf{7}
\]

\subsection*{(e) $E(X^{(2)})$}
\[
\mathbf{E(X^{(2)}) = \mu^{(2)} = \begin{bmatrix} 2 \\ 1 \end{bmatrix}}
\]

\subsection*{(f) $E(BX^{(2)})$}
\[
E(BX^{(2)}) = B E(X^{(2)}) = \begin{bmatrix} 1 & -2 \\ 2 & -1 \end{bmatrix} \begin{bmatrix} 2 \\ 1 \end{bmatrix} 
= \begin{bmatrix} 1(2) - 2(1) \\ 2(2) - 1(1) \end{bmatrix} = \mathbf{\begin{bmatrix} 0 \\ 3 \end{bmatrix}}
\]

\subsection*{(g) $Cov(X^{(2)})$}
\[
\mathbf{Cov(X^{(2)}) = \Sigma_{22} = \begin{bmatrix} 9 & -2 \\ -2 & 4 \end{bmatrix}}
\]

\subsection*{(h) $Cov(BX^{(2)})$}
\[
Cov(BX^{(2)}) = B \Sigma_{22} B' 
= \begin{bmatrix} 1 & -2 \\ 2 & -1 \end{bmatrix} \begin{bmatrix} 9 & -2 \\ -2 & 4 \end{bmatrix} \begin{bmatrix} 1 & 2 \\ -2 & -1 \end{bmatrix}
\]
\[
= \begin{bmatrix} 13 & -10 \\ 20 & -8 \end{bmatrix} \begin{bmatrix} 1 & 2 \\ -2 & -1 \end{bmatrix}
= \mathbf{\begin{bmatrix} 33 & 36 \\ 36 & 48 \end{bmatrix}}
\]

\subsection*{(i) $Cov(X^{(1)}, X^{(2)})$}
\[
\mathbf{Cov(X^{(1)}, X^{(2)}) = \Sigma_{12} = \begin{bmatrix} 2 & 2 \\ 1 & 0 \end{bmatrix}}
\]

\subsection*{(j) $Cov(AX^{(1)}, BX^{(2)})$}
\[
Cov(AX^{(1)}, BX^{(2)}) = A \Sigma_{12} B'
= \begin{bmatrix} 1 & 2 \end{bmatrix} \begin{bmatrix} 2 & 2 \\ 1 & 0 \end{bmatrix} \begin{bmatrix} 1 & 2 \\ -2 & -1 \end{bmatrix}
\]
\[
= \begin{bmatrix} 4 & 2 \end{bmatrix} \begin{bmatrix} 1 & 2 \\ -2 & -1 \end{bmatrix}
= \begin{bmatrix} 4(1) + 2(-2) & 4(2) + 2(-1) \end{bmatrix} = \mathbf{\begin{bmatrix} 0 & 6 \end{bmatrix}}
\]

% ---------------------------------------------------------
% PROBLEM 2.41
% ---------------------------------------------------------
\section*{Problem 2.41}
Given $\mu_X = \begin{bmatrix} 3 \\ 2 \\ -2 \\ 0 \end{bmatrix}$, $\Sigma_X = 3I_4 = \begin{bmatrix} 3 & 0 & 0 & 0 \\ 0 & 3 & 0 & 0 \\ 0 & 0 & 3 & 0 \\ 0 & 0 & 0 & 3 \end{bmatrix}$, and $A = \begin{bmatrix} 1 & -1 & 0 & 0 \\ 1 & 1 & -2 & 0 \\ 1 & 1 & 1 & -3 \end{bmatrix}$.

\subsection*{(a) Find $E(AX)$}
\[
E(AX) = A \mu_X = \begin{bmatrix} 1 & -1 & 0 & 0 \\ 1 & 1 & -2 & 0 \\ 1 & 1 & 1 & -3 \end{bmatrix} \begin{bmatrix} 3 \\ 2 \\ -2 \\ 0 \end{bmatrix}
\]
\[
= \begin{bmatrix} 1(3) - 1(2) + 0 + 0 \\ 1(3) + 1(2) - 2(-2) + 0 \\ 1(3) + 1(2) + 1(-2) - 3(0) \end{bmatrix}
= \mathbf{\begin{bmatrix} 1 \\ 9 \\ 3 \end{bmatrix}}
\]

\subsection*{(b) Find $Cov(AX)$}
\[
Cov(AX) = A \Sigma_X A' = A (3I_4) A' = 3 A A'
\]
First compute $AA'$:
\[
A A' = \begin{bmatrix} 1 & -1 & 0 & 0 \\ 1 & 1 & -2 & 0 \\ 1 & 1 & 1 & -3 \end{bmatrix} \begin{bmatrix} 1 & 1 & 1 \\ -1 & 1 & 1 \\ 0 & -2 & 1 \\ 0 & 0 & -3 \end{bmatrix}
= \begin{bmatrix} 2 & 0 & 0 \\ 0 & 6 & 0 \\ 0 & 0 & 12 \end{bmatrix}
\]
Multiplying by 3:
\[
Cov(AX) = 3 \begin{bmatrix} 2 & 0 & 0 \\ 0 & 6 & 0 \\ 0 & 0 & 12 \end{bmatrix} = \mathbf{\begin{bmatrix} 6 & 0 & 0 \\ 0 & 18 & 0 \\ 0 & 0 & 36 \end{bmatrix}}
\]

\subsection*{(c) Which pairs of linear combinations have zero covariances?}
Since $Cov(AX)$ is a diagonal matrix, all off-diagonal entries are zero. Therefore, \textbf{all pairs of linear combinations} have zero covariance:
\begin{itemize}
    \item Row 1 and Row 2: $Cov((AX)_1, (AX)_2) = 0$
    \item Row 1 and Row 3: $Cov((AX)_1, (AX)_3) = 0$
    \item Row 2 and Row 3: $Cov((AX)_2, (AX)_3) = 0$
\end{itemize}

\end{document}



